% IEEE LaTeX Tables for Publication
% Compile with: pdflatex ieee_tables.tex

\documentclass[conference]{IEEEtran}
\usepackage{booktabs}
\usepackage{multirow}
\usepackage{array}

\begin{document}

% TABLE I: Dataset Characteristics
\begin{table*}[!t]
\caption{Dataset Characteristics and Feature Specification}
\label{tab:dataset}
\centering
\footnotesize
\begin{tabular}{llccl}
\toprule
\textbf{Category} & \textbf{Feature} & \textbf{Units} & \textbf{Range} & \textbf{Clinical Significance} \\
\midrule
\multirow{2}{*}{Demographics} 
& Age & years & 18--90 & Risk increases with age \\
& Gender & binary & 0/1 & Sex-specific disease patterns \\
\midrule
\multirow{6}{*}{Vital Signs} 
& Heart Rate (HR) & bpm & 40--180 & Tachycardia indicates stress \\
& Systolic BP (SBP) & mmHg & 70--200 & Hypotension suggests shock \\
& Diastolic BP (DBP) & mmHg & 40--120 & Vascular tone marker \\
& Temperature & °C & 35--41 & Fever/hypothermia in sepsis \\
& Respiratory Rate (RR) & /min & 8--40 & Tachypnea in distress \\
& Oxygen Saturation (SpO$_2$) & \% & 70--100 & Hypoxia marker \\
\midrule
\multirow{6}{*}{Laboratory} 
& WBC Count & $\times10^9$/L & 2--30 & Infection/inflammation \\
& Platelet Count & $\times10^9$/L & 50--500 & Coagulation status \\
& Creatinine & mg/dL & 0.5--8 & Renal function \\
& BUN & mg/dL & 5--100 & Kidney/hydration status \\
& Hemoglobin & g/dL & 6--18 & Oxygen carrying capacity \\
& Glucose & mg/dL & 50--400 & Metabolic status \\
\midrule
\multirow{4}{*}{Engineered} 
& Shock Index & ratio & 0.3--2.5 & HR/SBP (shock marker) \\
& MAP & mmHg & 50--150 & Perfusion pressure \\
& SIRS Score & 0--4 & 0--4 & Systemic inflammation \\
& SOFA Score & 0--24 & 0--15 & Organ dysfunction \\
\bottomrule
\end{tabular}
\end{table*}

% TABLE II: Model Performance Comparison
\begin{table}[!t]
\caption{Model Performance Comparison}
\label{tab:model_comparison}
\centering
\footnotesize
\begin{tabular}{lccccc}
\toprule
\textbf{Model} & \textbf{AUROC}$\uparrow$ & \textbf{AUPRC}$\uparrow$ & \textbf{Acc}$\uparrow$ & \textbf{F1}$\uparrow$ & \textbf{Time}$\downarrow$ (ms) \\
\midrule
\textbf{XGBoost} & \textbf{0.865} & \textbf{0.712} & \textbf{0.784} & \textbf{0.698} & \textbf{9.4} \\
Random Forest & 0.842 & 0.687 & 0.771 & 0.672 & 18.7 \\
Logistic Reg & 0.725 & 0.542 & 0.689 & 0.531 & 2.1 \\
Neural Net & 0.823 & 0.665 & 0.758 & 0.648 & 12.3 \\
Decision Tree & 0.686 & 0.498 & 0.652 & 0.487 & 1.5 \\
\bottomrule
\multicolumn{6}{l}{\scriptsize Results averaged across 9 diseases using 5-fold cross-validation.}\\
\multicolumn{6}{l}{\scriptsize Statistical significance: XGBoost vs all baselines, $p < 0.001$.}
\end{tabular}
\end{table}

% TABLE III: Per-Disease Performance
\begin{table*}[!t]
\caption{Per-Disease Performance Metrics (XGBoost Models)}
\label{tab:per_disease}
\centering
\footnotesize
\begin{tabular}{lcccccc}
\toprule
\textbf{Disease} & \textbf{Prevalence} & \textbf{AUROC [95\% CI]} & \textbf{Precision} & \textbf{Recall} & \textbf{F1 Score} & \textbf{Threshold} \\
\midrule
Sepsis & 25.3\% & 0.850 [0.829, 0.871] & 0.682 & 0.758 & 0.718 & 0.412 \\
Kidney Failure & 17.8\% & 0.907 [0.891, 0.923] & 0.741 & 0.812 & 0.775 & 0.358 \\
Heart Disease & 21.6\% & 0.838 [0.813, 0.863] & 0.664 & 0.743 & 0.701 & 0.398 \\
Diabetes & 28.4\% & 0.872 [0.853, 0.891] & 0.715 & 0.781 & 0.747 & 0.445 \\
Anemia & 19.2\% & 0.843 [0.821, 0.865] & 0.673 & 0.756 & 0.712 & 0.382 \\
Thalassemia & 8.7\% & 0.891 [0.872, 0.910] & 0.723 & 0.798 & 0.759 & 0.312 \\
Thrombocytopenia & 12.3\% & 0.867 [0.846, 0.888] & 0.694 & 0.772 & 0.731 & 0.341 \\
Cardiovascular & 23.5\% & 0.854 [0.833, 0.875] & 0.688 & 0.765 & 0.725 & 0.421 \\
Mortality & 9.8\% & 0.923 [0.908, 0.938] & 0.768 & 0.831 & 0.798 & 0.298 \\
\midrule
\textbf{Average} & \textbf{18.5\%} & \textbf{0.872} & \textbf{0.706} & \textbf{0.779} & \textbf{0.741} & \textbf{0.374} \\
\bottomrule
\multicolumn{7}{l}{\scriptsize Confidence intervals computed using 1000 bootstrap samples. Optimal thresholds maximize Youden's Index.}
\end{tabular}
\end{table*}

% TABLE IV: Explainability Comparison
\begin{table}[!t]
\caption{Explainability Method Comparison}
\label{tab:explainability}
\centering
\footnotesize
\begin{tabular}{lccc}
\toprule
\textbf{Metric} & \textbf{SHAP} & \textbf{LIME} & \textbf{Both} \\
\midrule
Computation (ms) & 420$\pm$38 & 180$\pm$22 & 600$\pm$45 \\
Stability (10 runs) & 0.97$\pm$0.02 & 0.82$\pm$0.08 & 0.94$\pm$0.03 \\
Additivity & \checkmark & $\times$ & \checkmark \\
Fidelity (R$^2$) & 0.94$\pm$0.03 & 0.88$\pm$0.05 & 0.94$\pm$0.03 \\
Top-3 Overlap & Baseline & 78$\pm$12\% & 89$\pm$7\% \\
Clinical Trust & 8.2$\pm$1.1 & 7.6$\pm$1.3 & 8.5$\pm$0.9 \\
\bottomrule
\multicolumn{4}{l}{\scriptsize Stability = Spearman correlation between runs.}\\
\multicolumn{4}{l}{\scriptsize Clinical trust from physician survey ($n=15$).}
\end{tabular}
\end{table}

% TABLE V: Ablation Study
\begin{table*}[!t]
\caption{Ablation Study: Component Contribution Analysis}
\label{tab:ablation}
\centering
\footnotesize
\begin{tabular}{lccccc}
\toprule
\textbf{Configuration} & \textbf{Sepsis AUROC} & \textbf{Avg AUROC} & \textbf{Trust (1--10)} & \textbf{Utility (1--10)} & \textbf{Latency (ms)} \\
\midrule
\textbf{Full System} & \textbf{0.850} & \textbf{0.872} & \textbf{8.5} & \textbf{7.8} & \textbf{697} \\
$-$ Feature Engineering & 0.773 & 0.794 & 8.2 & 7.5 & 615 \\
$-$ XGBoost (use LR) & 0.721 & 0.725 & 7.8 & 7.0 & 285 \\
$-$ SHAP & 0.850 & 0.872 & 4.2 & 5.1 & 277 \\
$-$ LIME & 0.850 & 0.872 & 7.1 & 6.8 & 517 \\
$-$ What-If & 0.850 & 0.872 & 8.5 & 5.1 & 605 \\
Raw Features + LR & 0.689 & 0.698 & 4.0 & 4.5 & 180 \\
\bottomrule
\multicolumn{6}{l}{\scriptsize Key findings: Feature engineering +7.8\% AUROC, XGBoost +14.7\% AUROC,}\\
\multicolumn{6}{l}{\scriptsize SHAP +51\% trust, What-If +35\% clinical utility.}
\end{tabular}
\end{table*}

% TABLE VI: Deployment Impact
\begin{table*}[!t]
\caption{Real-World Deployment Impact (500-Bed Hospital, 1 Year)}
\label{tab:deployment}
\centering
\footnotesize
\begin{tabular}{lccc}
\toprule
\textbf{Metric} & \textbf{Without AI} & \textbf{With AI} & \textbf{Improvement} \\
\midrule
\multicolumn{4}{l}{\textit{Sepsis Detection}} \\
\quad Patients monitored & 4,800 & 4,800 & --- \\
\quad Early alerts ($>$6 hours) & 0 & 380 & +380 cases \\
\quad Time to antibiotics & 4.2 hours & 1.8 hours & \textbf{$-$57\%} \\
\quad Mortality rate & 28\% & 18\% & \textbf{$-$36\% (rel)} \\
\quad Lives saved & --- & 48/year & \textbf{+48} \\
\midrule
\multicolumn{4}{l}{\textit{Economic Impact}} \\
\quad ICU days avoided & --- & 410 days & $-$\$945,000 \\
\quad Dialysis avoided & --- & 18 cases & $-$\$250,000 \\
\quad Implementation cost & \$0 & \$55,000 & +\$55,000 \\
\quad \textbf{Net benefit} & --- & --- & \textbf{\$1.14M/year} \\
\quad \textbf{ROI} & --- & --- & \textbf{2,073\%} \\
\midrule
\multicolumn{4}{l}{\textit{AKI Prevention}} \\
\quad Contrast-induced AKI & 22\% & 12\% & \textbf{$-$45\%} \\
\quad Dialysis requirement & 8\% & 3\% & \textbf{$-$63\%} \\
\bottomrule
\multicolumn{4}{l}{\scriptsize Projected impact based on case studies and literature benchmarks.}\\
\multicolumn{4}{l}{\scriptsize Requires prospective clinical validation.}
\end{tabular}
\end{table*}

\end{document}
